\section{Accessible Ball-Union Channels}
\label{app:accessible}

\subsection{A geometric access guarantee}
The following construction gives an explicit passage through an object.
It does not assert a cooling rate, pressure drop, or manufacturing tolerance.
Let $e$ be a unit vector, let $m\geq1$ be an integer, and put
$c_j=c_0+jhe$ for $j=0,\ldots,m$, where $0<h<2r$. Define the open channel
$C=\bigcup_{j=0}^m B_r(c_j)$ and
$\rho=\sqrt{r^2-h^2/4}>0$.

\begin{proposition}[A positive-width passage]
\label{prop:channelcore}
The channel $C$ is path connected and contains the straight tube
\[
 T=\{c_0+te+v:0\leq t\leq mh,\ v\cdot e=0,\ |v|<\rho\}.
\]
If $c_0$ and $c_m$ lie strictly on opposite sides of a slab containing
the closed object, the removed set $C$ supplies an open route between
the two exterior sides. In particular, the route has positive width;
it is not merely a point-contact chain.
\end{proposition}
\begin{proof}
\textbf{Step 1: choose a neighboring center.}
For each $t\in[0,mh]$, rounding $t/h$ to a nearest integer gives
$j\in\{0,\ldots,m\}$ with $|t-jh|\leq h/2$.
\textbf{Step 2: bound the distance.}
Orthogonality of $v$ and $e$ gives
\[
 |c_0+te+v-c_j|^2=(t-jh)^2+|v|^2
 <h^2/4+\rho^2=r^2.
\]
Thus every point of $T$ belongs to an open ball in $C$.
\textbf{Step 3: connect all channel points.}
Any point of a ball can be connected to its center by a segment in that
ball. The centerline belongs to $T$. Concatenating two radial segments
with a centerline segment connects any two points of $C$.
\textbf{Step 4: reach the exterior.}
The endpoint centers lie in the exterior by the slab assumption.
Following the centerline from one endpoint to the other therefore
gives a route in $C$, including its portions removed from the object.
The strict tube inequality establishes positive transverse clearance.
Intersecting the tube with the object need not give a single component
on an arbitrary shape, but every such portion communicates with the
same open route. This statement does not require the remaining material
to be connected.
\end{proof}

\begin{proposition}[Protection and channel separation]
\label{prop:channelprotection}
Let $a>0$ and $g\geq0$. For parallel channels of radius $r$, a transverse distance at least
$r+a+g$ from a protected ball's center to the channel axis gives
a gap of at least $g$ from that radius-$a$ ball. An axis separation
of at least $2r+g$ gives a gap of at least $g$ between channels.
\end{proposition}
\begin{proof}
\textbf{Step 1: enclose the channel.}
Every channel ball lies in the radius-$r$ cylinder around its axis,
so the union lies in that cylinder.
\textbf{Step 2: project protected points.}
Orthogonal projection onto the plane perpendicular to the axis cannot
increase distances. A point in the protected ball projects at distance
at most $a$ from its center's projection. A channel point projects at
distance at most $r$ from its axis. The reverse triangle inequality
therefore bounds their separation below by $(r+a+g)-a-r=g$.
\textbf{Step 3: apply the same argument to two channels.}
Their two projected radial distances are each at most $r$, giving
the lower bound $(2r+g)-2r=g$.
\end{proof}

\subsection{What the construction does not prove}
For the nonnegative-source, zero-Dirichlet model and a protected
observation region, domain removal still gives a nonincreasing
regional objective, provided the weak solutions and trace comparison
used in \cref{app:regional} apply. To see this, set $D'=D\setminus
\overline C$ and use exactly the difference $u_D-u_{D'}$ and its
negative-part test from that proof. The new trace is nonnegative and
the source agrees on $D'$, so the weighted integral difference is
nonnegative. No assertion that the channel is infinitesimal is needed
for this comparison.

However, $-u(z)p(z)$ is the coefficient for a separated small ball.
A channel consists of many overlapping balls and meets the exterior.
The single-ball expansion therefore does not justify summing their
coefficients or assigning a certified finite benefit to the whole
channel. The experiment uses the point score as a placement heuristic,
then independently evaluates the actual channel geometry. A stronger
channel-shape derivative or finite-insertion theorem remains separate
work. Nor does an open route justify ideal isothermal walls: flow,
convective heat transfer, external manifolds, and pumping power are
absent from the model.

\subsection{Reproducible experiment}
The new comparison uses two sensor positions, five paired candidate
seeds, and five methods. Each design adds six vertical channels of
radius $0.075$. Eighteen balls span $z=-0.45$ to $z=0.825$ with
spacing $0.075$, so \cref{prop:channelcore} gives the exact core radius
$0.075\sqrt{3}/2$. Endpoint coordinates are checked against the
actual boundary mesh's bounds. Channel axes avoid the sensor ball
and three radius-$0.07$ component balls, with a $0.02$ additional
gap; channels have a $0.025$ mutual gap. Components are centered at
the three prescribed source peaks. The Gaussian source tails remain
a volumetric field on surviving material; this is not a fixed-power
electronic-component model.

Each seed gives 512 common candidate axes. Scores are queried at
$z=-0.03$ in the current domain and all channel balls are added after
each selection. Sparse \method{} and WoS hottest-point placement use
12 scout walks and 192 fresh rescore walks on 24 candidates.
Multigrid evaluates the same pointwise topological coefficient at
$57\times37\times29$ or $113\times73\times57$ nodes; random feasible
placement supplies a geometric control. Both solver families retain
the outer mesh through each complete six-round design. Timings include
candidate preparation, geometry setup, adaptive selection, I/O,
and process startup and shutdown. Reference evaluations are separate.

Four independently scrambled 128-point sensor-probe blocks and 256
fresh walks per probe evaluate every final design. A separate
$113\times73\times57$ grid provides another numerical reference.
First-seed designs additionally receive a fourfold shell refinement
and a $169\times109\times85$ grid check. Neither reference is exact.
Four independent 32,768-point scrambled blocks estimate removed
material volume and the removed fraction of the prescribed source.
All methods share channel count, radius, axial extent, and constraints;
the nonuniform housing thickness means their removed volumes are not
exactly equal. We report those volumes rather than describing this as
an equal-volume comparison. Port positions are optimized, not fixed
connections to an external manifold.

\begin{table*}[t]
\centering\small
\caption{Accessible-channel comparison. Regional temperature reductions are means $\pm$ one standard error across five designs per sensor. Time and removed material are averaged across both sensors. All methods use six channels with the same radius and extent; volume varies with local housing thickness. Point scores are placement proxies, not derivatives of an entire channel.}
\label{tab:accessible}
\begin{tabular}{lrrrr}
\toprule
Method & Left reduction (\%) & Right reduction (\%) & Design time (s) & Removed material (\%) \\
\midrule
Random & $17.30\pm2.23$ & $10.74\pm2.48$ & 0.26 & 4.28 \\
WoS hottest & $43.78\pm1.89$ & $30.39\pm4.17$ & 1.00 & 4.24 \\
TWoS proxy & $55.74\pm1.35$ & $52.47\pm0.56$ & 1.02 & 4.22 \\
MG proxy (57) & $58.14\pm1.11$ & $56.95\pm0.28$ & 1.26 & 4.37 \\
MG proxy (113) & $58.14\pm1.11$ & $57.35\pm0.41$ & 6.55 & 4.36 \\
\bottomrule
\end{tabular}
\end{table*}

The saved-data audit verifies all 50 channel geometries and 240 nonrandom selections. All walks terminate before the step cap and all recorded multigrid solves meet the $10^{-8}$ residual tolerance. The largest first-seed shell-refinement change is $0.084$ percentage points; the grid-refinement change is $1.721$ points. Across all designs, the two primary numerical references differ by at most $3.133$ points. These are sensitivity checks, not bounds on the remaining numerical error.

